\section{Gauss points over arbitrary seminorm spectra}\label{sec:general-gauss-points}

Let $A$ be any commutative ring, $y\in\mathcal S(A)$, and $r\in\R_{\ge0}$.
All constructions include $r=0$ and require no field or finite type hypothesis.
For relative statements, $A$ is an algebra over a commutative ring $R$ with
a prescribed real-valued ultrametric seminorm $v$.

\begin{lemma}[Multiplicativity of the Gauss seminorm]\label{spectrum:gauss-valuation}
\lean{EnvelopeProperty.gaussValuation, EnvelopeProperty.gaussFun_mul}
\leanok
\uses{spectrum:affine-seminorm}
For a real-valued multiplicative ultrametric seminorm $w$ on $A$, the formula
\[
 N_r(F)=\max_i w(F_i)r^i,\qquad F=\sum_iF_iX^i\in A[X],
\]
defines a multiplicative ultrametric seminorm on $A[X]$.
\end{lemma}
\begin{proof}
The values at zero and one follow from the formula. Applying the ultrametric
inequality to each coefficient proves the ultrametric inequality for $N_r$.
For each degree $n$, the coefficient of $FG$ is $\sum_{a+b=n}F_aG_b$;
its weighted value is at most $\max_{a+b=n}w(F_a)r^a w(G_b)r^b$
and thus at most $N_r(F)N_r(G)$. This proves the product upper bound.
If either norm is zero, that upper bound proves multiplicativity immediately.
Otherwise choose the least indices $i,j$ attaining the respective positive maxima.
The term $F_iG_j$ in the coefficient of degree $i+j$ has weighted value
$N_r(F)N_r(G)>0$, so $r^{i+j}>0$.
For any other pair $a+b=i+j$, if $a<i$ then $w(F_a)r^a<N_r(F)$;
if $a>i$ then $b<j$ and $w(G_b)r^b<N_r(G)$.
Since the other factor is bounded by its positive maximum, each such term
has strictly smaller weighted value. Dividing the comparison by the positive
$r^{i+j}$ shows that $F_iG_j$ is the unique term of greatest seminorm in that
coefficient. The ultrametric inequality, applied also after subtracting the
other terms, implies that the seminorm of the coefficient equals
$w(F_iG_j)$. Its weighted value is therefore $N_r(F)N_r(G)$.
This gives the reverse product bound and proves multiplicativity.
\end{proof}

\begin{definition}[Projection and Gauss points]\label{spectrum:gauss-point}
\lean{EnvelopeProperty.Berkovich.Spectrum.proj,
EnvelopeProperty.Berkovich.Spectrum.gauss,
EnvelopeProperty.Berkovich.RelativeSpectrum.proj,
EnvelopeProperty.Berkovich.RelativeSpectrum.gauss}
\leanok
\uses{spectrum:gauss-valuation,spectrum:pullback,spectrum:relative}
Restriction along $A\to A[X]$ gives a projection $p:\mathcal S(A[X])\to\mathcal S(A)$.
The Gauss point $\eta_{y,r}$ has seminorm $|F|_{\eta_{y,r}}=\max_i|F_i|_y r^i$.
It satisfies $p(\eta_{y,r})=y$ and $|X|_{\eta_{y,r}}=r$.
If $y\in\mathcal S_v(A)$, then $\eta_{y,r}\in\mathcal S_v(A[X])$ with the same values
on the prescribed base ring.
\end{definition}
\begin{proof}
The preceding lemma defines the Gauss seminorm. The value of a constant polynomial
$a$ is $|a|_y$ and the value of $X$ is $r$, proving the stated identities.
In particular, the values of constants coming from $R$ remain $v(r')$ for $r'\in R$,
so the construction preserves the relative locus.
\end{proof}

\begin{lemma}[The zero-radius Gauss point]\label{spectrum:gauss-zero-radius}
\lean{EnvelopeProperty.Berkovich.Spectrum.gauss_zero_eq_comap,
EnvelopeProperty.Berkovich.RelativeSpectrum.gauss_zero_eq_comap}
\leanok
\uses{spectrum:gauss-point}
At radius zero, $\eta_{y,0}$ is pullback of $y$ along evaluation $A[X]\to A$, $X\mapsto0$.
\end{lemma}
\begin{proof}
In the Gauss maximum at zero, the degree-zero term is $|F_0|_y$ and every
positive-degree term is zero. Thus $|F|_{\eta_{y,0}}=|F(0)|_y$ for every $F$,
which identifies the two seminorms. This identity also holds on each relative locus.
\end{proof}

\begin{lemma}[Maximality on a fibre disc]\label{spectrum:gauss-maximality}
\lean{EnvelopeProperty.Berkovich.Spectrum.val_le_gauss_val,
EnvelopeProperty.Berkovich.RelativeSpectrum.val_le_gauss_val}
\leanok
\uses{spectrum:gauss-point}
If $p(z)=y$ and $|X|_z\le r$, then $|F|_z\le|F|_{\eta_{y,r}}$ for every polynomial $F$.
This applies in both the absolute and relative spectra.
\end{lemma}
\begin{proof}
Write $F=\sum_iF_iX^i$. The ultrametric inequality bounds its seminorm at $z$
by $\max_i|F_i|_z|X|_z^i$. Since $p(z)=y$, the constant coefficient evaluations
at $z$ equal those at $y$, and $|X|_z^i\le r^i$.
The bound is consequently at most $\max_i|F_i|_y r^i$, as asserted.
\end{proof}

\begin{lemma}[Joint continuity]\label{spectrum:gauss-joint-continuity}
\lean{EnvelopeProperty.Berkovich.Spectrum.continuous_gauss_prod,
EnvelopeProperty.Berkovich.RelativeSpectrum.continuous_gauss_prod}
\leanok
\uses{spectrum:gauss-point,spectrum:topology}
The map $(y,r)\mapsto\eta_{y,r}$ is jointly continuous, including at radius zero,
on both $\mathcal S(A)\times\R_{\ge0}$ and $\mathcal S_v(A)\times\R_{\ge0}$.
\end{lemma}
\begin{proof}
For each fixed polynomial, its Gauss evaluation is the finite maximum of the
functions $(y,r)\mapsto |F_i|_y r^i$. Each is continuous by continuity of
evaluation, multiplication, and nonnegative integer powers. A finite maximum
is continuous. The evaluation characterization of the spectrum topology proves
joint continuity. The relative topology is the induced subspace topology, so
the same argument proves the relative statement.
\end{proof}

\begin{theorem}[Compact closed fibre discs]\label{spectrum:gauss-compact-fibre-disc}
\lean{EnvelopeProperty.Berkovich.Spectrum.isCompact_fibreDisc,
EnvelopeProperty.Berkovich.RelativeSpectrum.isCompact_fibreDisc}
\leanok
\uses{spectrum:gauss-maximality,spectrum:compact-coordinate-bounds,spectrum:compact-relative-bounds}
For every $y$ and every $r\ge0$, the closed fibre disc
$\{z:p(z)=y,\ |X|_z\le r\}$ is compact in either spectrum.
\end{theorem}
\begin{proof}
The fibre condition is closed because the projection is continuous and the
base spectrum is Hausdorff. The radius condition is closed by continuity of
evaluation at $X$. Maximality bounds the evaluation of every polynomial $F$
on this disc by the finite number $|F|_{\eta_{y,r}}$.
The compact coordinate-bound theorem therefore places the disc inside a
compact set; its closedness makes it compact. In the relative case use the
relative coordinate-bound theorem and the relative projection.
\end{proof}
